\SourceLeaf{038}
\Article{27}
The function \(\Omega\) can be expressed in indefinitely many forms, which it will be worth developing. First, by squaring equations I of art. 20 and adding, we obtain immediately
\[
\begin{aligned}
\Omega={}&xx\Sigma aa+yy\Sigma bb+zz\Sigma cc+\text{etc.}\\
&+2xy\Sigma ab+2xz\Sigma ac+2yz\Sigma bc+\text{etc.}\\
&+2x\Sigma al+2y\Sigma bl+2z\Sigma cl+\text{etc.}+\Sigma ll
\end{aligned}
\]
which is the \textit{first} form.

Multiplying the same equations respectively by \(v,v',v''\) etc. and adding, we obtain:
\[
\Omega=\xi x+\eta y+\zeta z+\text{etc.}+lv+l'v'+l''v''+\text{etc.}
\]
and hence, substituting for \(v,v',v''\) etc. the expressions given in the preceding article,
\[
\Omega=\xi x+\eta y+\zeta z+\text{etc.}-A\xi-B\eta-C\zeta-\text{etc.}+M
\]
or
\[
\Omega=\xi(x-A)+\eta(y-B)+\zeta(z-C)+\text{etc.}+M
\]
which is the \textit{second} form.

Substituting expressions VII of art. 21 for \(x-A,y-B,z-C\) etc. in the second form, we obtain the \textit{third} form:
\[
\begin{aligned}
\Omega={}&[\alpha\alpha]\xi\xi+[\beta\beta]\eta\eta+[\gamma\gamma]\zeta\zeta+\text{etc.}\\
&+2[\alpha\beta]\xi\eta+2[\alpha\gamma]\xi\zeta+2[\beta\gamma]\eta\zeta+\text{etc.}+M
\end{aligned}
\]
To these we may add the \textit{fourth} form, which follows of itself from the third form and the formulas of the preceding article:
\[
\Omega=(v-\lambda)^2+(v'-\lambda')^2+(v''-\lambda'')^2+\text{etc.}+M,\quad\text{or}
\]
\[
\Omega=M+\Sigma(v-\lambda)^2
\]
This form puts the condition for a minimum directly before our eyes.
\Article{28}
Let \(e,e',e''\) etc. be the errors committed in the observations that gave \(V=L,\allowbreak V'=L'\quad V''=L''\) etc.; that is, let the true values of the functions \(V,V',V''\) etc. be respectively \(L-e,L'-e',L''-e''\) etc., and hence the true values of \(v,v',v''\) etc. respectively \(-e\BareRoot{p},\allowbreak-e'\BareRoot{p'},\allowbreak-e''\BareRoot{p''}\) etc. The true value of \(x\) will therefore be
\SourceLeaf{039}
\[
=A-\alpha e\BareRoot{p}-\alpha'e'\BareRoot{p'}-\alpha''e''\BareRoot{p''}-\text{etc.}
\]
or, for the error committed in the best determination of the value of \(x\), which it is convenient to denote by \(Ex\), we have
\[
=\alpha e\BareRoot{p}+\alpha'e'\BareRoot{p'}+\alpha''e''\BareRoot{p''}+\text{etc.}
\]
Similarly, the error committed in the best determination of the value of \(y\), which we shall denote by \(Ey\), will be
\[
=\beta e\BareRoot{p}+\beta'e'\BareRoot{p'}+\beta''e''\BareRoot{p''}+\text{etc.}
\]
The mean value of the square \((Ex)^2\) is found to be
\[
=mmp(\alpha\alpha+\alpha'\alpha'+\alpha''\alpha''+\text{etc.})=mmp[\alpha\alpha]
\]
and for the mean value of the square \((Ey)^2\) we similarly find \(=mmp[\beta\beta]\) etc., as we have already shown above. We may now also assign the mean value of the \textit{product} \(Ex.Ey\), which is found to be
\[
=mmp(\alpha\beta+\alpha'\beta'+\alpha''\beta''+\text{etc.})=mmp[\alpha\beta]
\]
These statements can also be expressed neatly as follows. The mean values of the squares \((Ex)^2,(Ey)^2\) etc. are respectively equal to the products of \(\frac12mmp\) and the second partial derivatives
\[
\frac{dd\Omega}{d\xi^2},\quad\frac{dd\Omega}{d\eta^2}\quad\text{etc.}
\]
and the mean value of a product such as \(Ex.Ey\) is equal to the product of \(\frac12mmp\) and the derivative \(\frac{dd\Omega}{d\xi.d\eta}\): provided, of course, that \(\Omega\) is considered as a function of the indeterminates \(\xi,\eta,\zeta\) etc.
\clearpage
